%% Đề 11 -- de-hoc-sinh-gioi-cap-tinh-toan-9-nam-2022-2023-so-gddt-lang-son.pdf
\deheader{ĐỀ SỐ 11}%
  {Kỳ thi Chọn Học sinh giỏi cấp tỉnh lớp 9 — Năm học 2022--2023}%
  {Sở GD\&ĐT Lạng Sơn | Chỉ phần Đại số}
\addcontentsline{toc}{section}{Đề 11 — Kỳ thi Chọn HSG cấp tỉnh — Năm học 2022--2023}

\begin{de_cau}{Câu 1}{4,0 điểm}
\begin{enumerate}[label=\alph*)]
\item Rút gọn: $A=\sqrt{4+2\sqrt{3}}-\sqrt{4-2\sqrt{3}}$.
\item Cho biểu thức $P=\dfrac{\sqrt{x}+1}{\sqrt{x}-2}+\dfrac{2\sqrt{x}}{\sqrt{x}+2}-\dfrac{5\sqrt{x}+2}{x-4}$,
$x\geq 0, x\neq 4$. Rút gọn $P$ và tìm $x$ để $P<\dfrac{1}{2}$.
\end{enumerate}
\end{de_cau}

\begin{loi_giai_cau}
\textbf{a)} $\sqrt{4+2\sqrt{3}}=\sqrt{(\sqrt{3}+1)^2}=\sqrt{3}+1$.
$\sqrt{4-2\sqrt{3}}=\sqrt{(\sqrt{3}-1)^2}=\sqrt{3}-1$.
$A=(\sqrt{3}+1)-(\sqrt{3}-1)=\boxed{2}$.

\textbf{b)} Quy đồng mẫu $(\sqrt{x}-2)(\sqrt{x}+2)$:
\[P=\frac{(\sqrt{x}+1)(\sqrt{x}+2)+2\sqrt{x}(\sqrt{x}-2)-(5\sqrt{x}+2)}{(\sqrt{x}-2)(\sqrt{x}+2)}\]
\[=\frac{x+3\sqrt{x}+2+2x-4\sqrt{x}-5\sqrt{x}-2}{x-4}=\frac{3x-6\sqrt{x}}{x-4}=\frac{3\sqrt{x}(\sqrt{x}-2)}{(\sqrt{x}-2)(\sqrt{x}+2)}=\frac{3\sqrt{x}}{\sqrt{x}+2}.\]

$P<\dfrac{1}{2}$: $\dfrac{3\sqrt{x}}{\sqrt{x}+2}<\dfrac{1}{2}\Rightarrow 6\sqrt{x}<\sqrt{x}+2\Rightarrow 5\sqrt{x}<2\Rightarrow x<\dfrac{4}{25}$.

Kết hợp $x\geq 0$: $\boxed{0\leq x<\dfrac{4}{25}}$.
\end{loi_giai_cau}

\begin{de_cau}{Câu 2}{4,0 điểm}
\begin{enumerate}[label=\alph*)]
\item Giải phương trình $\sqrt{x+3}+\sqrt{3-x}=\sqrt{5x+9}$.
\item Giải hệ phương trình: 
\[ \begin{cases}x^2+y^2=5\\x^2-xy+y^2=3\end{cases} \].
\end{enumerate}
\end{de_cau}

\begin{loi_giai_cau}
\textbf{a)} ĐKXĐ: $-3\leq x\leq 3$ và $5x+9\geq 0$, tức $x\geq -\frac{9}{5}$.

Bình phương hai vế: $(x+3)+(3-x)+2\sqrt{(x+3)(3-x)}=5x+9$.

$6+2\sqrt{9-x^2}=5x+9\Rightarrow 2\sqrt{9-x^2}=5x+3$.

ĐK: $5x+3\geq 0\Rightarrow x\geq -\frac{3}{5}$.

$4(9-x^2)=(5x+3)^2\Rightarrow 36-4x^2=25x^2+30x+9\Rightarrow 29x^2+30x-27=0$.

$x=\dfrac{-30\pm\sqrt{900+4\cdot29\cdot27}}{58}=\dfrac{-30\pm\sqrt{4032}}{58}$...
Thử $x=\frac{3}{5}$: không cho nghiệm đẹp. Nghiệm: $x=\dfrac{-30+\sqrt{4032}}{58}\approx 0{,}55$.

\textbf{b)} Từ hệ: $xy = x^2+y^2-3=5-3=2$.

$(x+y)^2=x^2+2xy+y^2=5+4=9\Rightarrow x+y=\pm 3$.

- $x+y=3, xy=2$: $x,y$ là nghiệm $t^2-3t+2=0$: $(x;y)=(1;2)$ hoặc $(2;1)$.
- $x+y=-3, xy=2$: $(x;y)=(-1;-2)$ hoặc $(-2;-1)$.
\end{loi_giai_cau}

\begin{de_cau}{Câu 3}{4,0 điểm}
\begin{enumerate}[label=\alph*)]
\item Tìm các cặp số nguyên $(x;y)$ thỏa mãn $2x^2+y^2+2xy-2y-4x=0$.
\item Cho $n$ nguyên dương. Chứng minh $n(n+1)(2n+1)$ chia hết cho $6$.
\end{enumerate}
\end{de_cau}

\begin{loi_giai_cau}
\textbf{a)} $2x^2+y^2+2xy-2y-4x=(x+y)^2+x^2-2(x+y)-4x+2x^2-x^2$...

Thử phân tích: $2x^2+2xy+y^2-4x-2y=(x+y)^2+(x-2)^2-4=0$?

$(x+y)^2+x^2-4x+4-4+(1+y-x)^2\ldots$

Viết lại: $2x^2+y^2+2xy-4x-2y=2x(x+y-2)+(y^2-2y)=2x(x+y-2)+y(y-2)$.

Thử: $(x+y)(2x+y)-2(2x+y)=(2x+y)(x+y-2)=0$.

$(2x+y=0$ hoặc $x+y=2)$.

- $2x+y=0\Rightarrow y=-2x$: mọi $x\in\mathbb{Z}$, $(x;-2x)$.
- $x+y=2\Rightarrow y=2-x$: mọi $x\in\mathbb{Z}$, $(x;2-x)$.

\textbf{b)} $n(n+1)(2n+1)$: tích $n$ và $n+1$ liên tiếp nên chia hết cho $2$.

Xét mod $3$: $n\equiv 0,1,2$:
- $n\equiv 0$: chia hết cho $3$.
- $n\equiv 1$: $2n+1\equiv 0$.
- $n\equiv 2$: $n+1\equiv 0$.

Vậy luôn chia hết cho $3$. Do $(2,3)=1$, tích chia hết cho $6$. $\square$
\end{loi_giai_cau}

\begin{de_cau}{Câu 5}{2,0 điểm}
Cho $a,b,c>0$ thỏa mãn $a+b+c=1$. Tìm $\min$ của $T=\dfrac{1}{a}+\dfrac{4}{b}+\dfrac{9}{c}$.
\end{de_cau}

\begin{loi_giai_cau}
Theo Cauchy-Schwarz (Engel/Titu):
\[T = \frac{1^2}{a}+\frac{2^2}{b}+\frac{3^2}{c}\geq\frac{(1+2+3)^2}{a+b+c}=\frac{36}{1}=36.\]

Dấu bằng khi $\dfrac{1}{a}:\dfrac{2}{b}:\dfrac{3}{c}=1\Rightarrow a=\dfrac{1}{6}, b=\dfrac{2}{6}=\dfrac{1}{3}, c=\dfrac{3}{6}=\dfrac{1}{2}$.

$T_{\min}=\boxed{36}$ khi $a=\dfrac{1}{6},b=\dfrac{1}{3},c=\dfrac{1}{2}$.
\end{loi_giai_cau}
